\title[Causal Evaluation of Decomposition Ensemble Stock Index Forecasting]{Causal Evaluation of Decomposition Ensemble Stock Index Forecasting: Reported Accuracy Depends on the Decomposition Seeing the Future}

%% AUTHORS: copied from the original Group 17 submission as a placeholder.
%% Confirm the list and order, add the members who carried out this study,
%% and supply every e-mail address (the \email lines marked TODO) before
%% submission. No address below has been invented

\author[1]{\fnm{Prayag} \sur{Sharma}}
\email{prayag.sharma2024@vitstudent.ac.in}

\author[1]{\fnm{Veer Vikram} \sur{Saxena}}
\email{veer.vikramsaxena2024@vitstudent.ac.in}

\author[1]{\fnm{Ayush} \sur{Srivastava}}
\email{ayush.srivastava2024@vitstudent.ac.in}

\author*[2]{\fnm{Ravi} \sur{Tiwari}}
\email{ravi.tiwari@vit.ac.in}

\affil[1]{%
\orgdiv{School of Computer Science and Engineering (SCOPE)},
\orgname{Vellore Institute of Technology (VIT)},
\orgaddress{\city{Chennai}, \state{Tamil Nadu}, \country{India}}%
}

\affil[2]{%
\orgdiv{School of Electronics Engineering (SENSE)},
\orgname{Vellore Institute of Technology (VIT)},
\orgaddress{\city{Chennai}, \state{Tamil Nadu}, \country{India}}%
}

\abstract{Decomposition ensemble models, which split a price series into
components with empirical or variational mode decomposition and forecast each
component with a neural network, report daily index forecasts with percentage
errors well below one tenth of one percent. A decomposition is not causal,
however: the value of a component on a given day depends on the days that
follow it, so a series decomposed in full before the training and test split
carries the test prices into the test inputs. We measure how much of the
reported accuracy depends on this. Reproducing a published model that combines
improved complete ensemble empirical mode decomposition with adaptive noise, a
second variational mode decomposition stage tuned by particle swarm
optimisation, and a bidirectional long short term memory network with self
attention feeding either a temporal convolutional head or a tiny recursive
head, we compare full series decomposition with a walk forward protocol that
redecomposes the available history at every forecast origin, on the daily
highest and lowest prices of the S\&P~500 and the Shanghai Stock Exchange
Composite Index over ten years. Before any network is trained, a linear
regression on the full series inputs explains 51 to 60 percent of the next
day's price change, and none of it under walk forward decomposition. Trained
on identical code, seeds and budgets, the networks reach a mean absolute
percentage error of 0.07 to 0.13 percent and a coefficient of determination
above 0.999 under full series decomposition and beat persistence on every
series by the modified Diebold and Mariano test; under walk forward
decomposition their error is 4.9 to 10.4 times larger and they fall
significantly behind persistence on every series. The reported advantage of
the convolutional head over the recursive head holds only under full series
decomposition. The evidence covers two indices, a one day horizon and one seed
per network. In reproducing the pipeline we identified and corrected twelve
implementation issues, which we document, and we propose two inexpensive
checks, an input diagnostic and a leakage ratio, that make the dependence of a
headline result on its protocol visible.}

\keywords{Deep learning, Financial time series, Look ahead bias, Boundary issue, Walk forward validation, Diebold and Mariano test}
